\subsection{Surface-informed polymer layer}
\label{sec:surfaceInformedPolymerLayer}

This verification folder exercises the surface-informed polymer as both a continuum material and a finite-thickness cohesive-zone material.  Four variants are run from one input: continuum tension, continuum compression, cohesive-zone tension, and cohesive-zone compression.  The continuum variants use a single-material polymer patch with 12 cells by 12 cells and a height equal to the cohesive-zone film thickness.  The deformation-gradient table is applied directly to that patch, avoiding artificial stiff-block/polymer interfaces and associated mixed-cell weak-discontinuity errors.  The cohesive-zone variants retain two stiff elastic blocks separated by a finite-thickness cohesive interface with the same film thickness.  The imposed cosine ramp is intentionally slow compared with the layer acoustic-transit time so the history is a quasi-static constitutive check rather than a wave-propagation benchmark.

The comparison curves use the continuum J2 radial-return update with retained mean stress for continuum variants and the reduced cohesive thin-film update with the same retained volumetric normal stress for cohesive-zone variants.  In the cohesive projection, the mean film stress $p=K\epsilon_n$ is retained and only the deviatoric normal and shear stresses are returned to the scalar flow surface.  The post-processor plots loading-direction stress in MPa versus reconstructed film strain and equivalent plastic strain versus reconstructed film strain.  For measured histories, the film-strain coordinate is reconstructed from the actual deformation-gradient or length history written by GEOS when available, with the input cosine loading table used only as a fallback.  When VisIt is available, it also renders loading-direction stress and equivalent plastic strain at four output states for each variant.

\IfFileExists{\CaseOutputDir/surface_polymer_layer_results.tex}{\input{\CaseOutputDir/surface_polymer_layer_results.tex}}{\paragraph{Generated results.} Results are produced by \texttt{postProcess\_surfaceInformedPolymerLayer.py} after the GEOS jobs finish.}
